\documentclass[12pt]{article}
\usepackage[utf8]{inputenc}
\usepackage[T1]{fontenc}
\usepackage{hyperref}
\title{Talossan Language Resources for K-12 Curriculum Development}
\author{Compiled by Claude Sonnet 5 (claude-sonnet-5) for Andr\'as Kornai}
\date{}
\begin{document}
\maketitle
\section{Introduction}
Talossan (ISO 639-3 tzl) is a constructed language created by Robert Ben Madison in 1980 for the micronation Kingdom of Talossa. It has no territorial community of native speakers. Searches found a full dictionary of more than 28,000 words and a small, active community of users. No school curriculum exists. Everything below comes from the community's own wiki and from general reference pages.
\section{Primary Educational Resources}
No K-12 textbooks or primers were found in search results. The Talossa wiki says a textbook and a dictionary appeared in 1987, but the searches did not return a copy of either. The Omniglot phrase page is the only learner-oriented item found: \url{https://omniglot.com/language/phrases/talossan.htm}.
\section{Linguistic and Reference Materials}
The Talossa wiki language page describes the language, its versions and its grammar: \url{https://wiki.talossa.com/Language}. It names a grammar, \emph{Dals Recomendăs sür el Glheþ Talossan}, by Marcel Eðo Pairescu Tafial and SIGN. Talossan 4.0 is reported to be in beta. The Linguifex entry gives an overview: \url{https://linguifex.com/wiki/Talossan}. The Wikipedia article on Talossa covers the micronation and its language: \url{https://EN.wikipedia.com/wiki/Talossa}.
\section{Religious and Community Resources}
The Comit\`a per l'\'Utzil del Gle\th{} (CUG), founded in 1983, issued decrees and vocabulary supplements and is now dormant: \url{https://wiki.talossa.com/CUG}. The Societ\`a per l'Ilesnaziun del Gle\th{} Naziunal (SIGN), set up in 2017, now leads the study of the language. No religious texts were found.
\section{Digital and Audio Resources}
No audio or video resources were identified in the search results.
\section{Curriculum Development Context}
Talossan is an artificial language used by a small community. A K-12 curriculum would have to be built from scratch on the dictionary and grammar kept by the community. Contact SIGN through the wiki to ask which materials are current. The sources found are thin.
\end{document}